\documentclass[11pt,a4paper]{article}

% ─────────────────────────────────────────────────────────────────────────
% Packages
% ─────────────────────────────────────────────────────────────────────────
\usepackage[utf8]{inputenc}
\usepackage[T1]{fontenc}
\usepackage{lmodern}
\usepackage{microtype}
\usepackage{geometry}
\geometry{margin=2.5cm}
\usepackage{amsmath,amssymb,amsthm}
\usepackage{mathtools}
\usepackage{algorithm}
\usepackage{algorithmic}
\usepackage{graphicx}
\usepackage{xcolor}
\usepackage{booktabs}
\usepackage{longtable}
\usepackage{array}
\usepackage{multirow}
\usepackage{float}
\usepackage{caption}
\usepackage{subcaption}
\usepackage{enumitem}
\usepackage{tikz}
\usetikzlibrary{shapes.geometric, arrows.meta, positioning, calc, fit, backgrounds}
\usepackage{hyperref}
\usepackage{cleveref}
\usepackage{fancyhdr}
\usepackage{titlesec}
\usepackage{setspace}

% ─────────────────────────────────────────────────────────────────────────
% Color Definitions
% ─────────────────────────────────────────────────────────────────────────
\definecolor{primaryblue}{RGB}{45, 62, 80}
\definecolor{accentgreen}{RGB}{39, 174, 96}
\definecolor{warnorange}{RGB}{230, 126, 34}
\definecolor{critred}{RGB}{192, 57, 43}
\definecolor{lightgray}{RGB}{236, 240, 241}

% ─────────────────────────────────────────────────────────────────────────
% Theorem Environments
% ─────────────────────────────────────────────────────────────────────────
\theoremstyle{definition}
\newtheorem{definition}{Definition}[section]
\newtheorem{theorem}{Theorem}[section]
\newtheorem{lemma}{Lemma}[section]
\newtheorem{property}{Property}[section]
\newtheorem{remark}{Remark}[section]

% ─────────────────────────────────────────────────────────────────────────
% Custom Commands
% ─────────────────────────────────────────────────────────────────────────
\newcommand{\tool}[1]{\texttt{#1}}
\newcommand{\code}[1]{\texttt{#1}}
\newcommand{\phantom}{\textsc{Phantom}}

% ─────────────────────────────────────────────────────────────────────────
% Document Information
% ─────────────────────────────────────────────────────────────────────────
\title{\textbf{LLM-Driven Orchestration for Automated Vulnerability Detection}\\\vspace{0.3em}
\large A ReAct-Based Architecture with Structured State Management}
\author{System Documentation}
\date{\today}

% ─────────────────────────────────────────────────────────────────────────
% Document Start
% ─────────────────────────────────────────────────────────────────────────
\begin{document}
\maketitle
\thispagestyle{empty}

\vspace{2em}
\begin{abstract}
\noindent
This thesis presents \phantom{}, an LLM-driven orchestration system for automated vulnerability detection. The system implements a bounded ReAct execution loop with structured state management, hypothesis tracking through exact-match deduplication, and dual-anchor memory preservation. We formalize the CVE Template Gap as a structural impossibility theorem for signature-based scanners and position \phantom{} as a reasoning-based alternative. The architecture is verified against the implementation codebase. Evaluation methodology is proposed but not executed.
\end{abstract}

\tableofcontents
\newpage

% ═════════════════════════════════════════════════════════════════════════
% CHAPTER 1: Introduction
% ═════════════════════════════════════════════════════════════════════════
\section{Introduction}
\label{sec:introduction}

\subsection{Context and Motivation}

Automated vulnerability detection relies on two distinct paradigms: signature-based scanners (e.g., Nuclei, Nessus) and manual penetration testing. Signature-based tools offer deterministic, reproducible results but are bounded by their template libraries. Manual testing provides semantic reasoning but does not scale.

Large Language Models (LLMs) present an intermediate approach: automated execution with reasoning capabilities. This thesis documents \phantom{}, a system exploring this design space.

\subsection{Research Questions}

\begin{enumerate}[label=\textbf{RQ\arabic*:}, leftmargin=2em]
    \item How can LLM reasoning be integrated with structured vulnerability tracking?
    \item What memory architecture preserves critical security findings across long-running scans?
    \item What are the fundamental limits of template-based vs. reasoning-based detection?
\end{enumerate}

\subsection{Contributions}

\begin{enumerate}[label=\textbf{C\arabic*:}, leftmargin=2em]
    \item \textbf{CVE Template Gap:} Formal definition of the structural impossibility of signature-based detection for non-templated vulnerabilities.
    \item \textbf{Dual-Anchor Memory:} Verified architecture with message-level and state-level anchor preservation.
    \item \textbf{Bounded ReAct Loop:} Implementation with explicit termination conditions and iteration limits.
    \item \textbf{Exact-Match Deduplication:} Hypothesis tracking with formal deduplication semantics.
\end{enumerate}

\subsection{Thesis Structure}

Section~\ref{sec:problem} defines the problem formally. Section~\ref{sec:sota} reviews related work. Section~\ref{sec:gaps} identifies research gaps. Section~\ref{sec:approach} presents the proposed approach. Section~\ref{sec:formal} provides formal models. Section~\ref{sec:architecture} details system architecture. Section~\ref{sec:design} discusses design decisions. Section~\ref{sec:limitations} documents limitations. Section~\ref{sec:discussion} provides discussion. Section~\ref{sec:conclusion} concludes.

% ═════════════════════════════════════════════════════════════════════════
% CHAPTER 2: Problem Definition
% ═════════════════════════════════════════════════════════════════════════
\section{Problem Definition}
\label{sec:problem}

\subsection{The Detection Problem}

Let $S$ be the set of all software systems. Let $V(s)$ be the set of vulnerabilities in system $s \in S$. The detection problem is:

\begin{definition}[Vulnerability Detection]
Given $s \in S$, compute a set $\hat{V}(s) \subseteq V(s)$ such that:
\begin{enumerate}[label=(\roman*)]
    \item Completeness: $\hat{V}(s)$ approximates $V(s)$ with bounded error
    \item Soundness: $\hat{V}(s) \subseteq V(s)$ (no false positives)
\end{enumerate}
\end{definition}

\subsection{Signature-Based Approach}

Let $T$ be a finite set of detection templates (signatures). The signature-based detector $D_T$ is:

\begin{equation}
D_T(s) = \{v \in V(s) : \exists t \in T, \text{matches}(t, v)\}
\end{equation}

\begin{property}[Bounded Detection]
For any $T$, $|D_T(s)| \leq |T|$.
\end{property}

\subsection{Reasoning-Based Approach}

Let $\mathcal{R}$ be a reasoning function. The reasoning-based detector $D_\mathcal{R}$ is:

\begin{equation}
D_\mathcal{R}(s) = \{v \in V(s) : \mathcal{R}(s, v) = \text{vulnerable}\}
\end{equation}

The reasoning function $\mathcal{R}$ is not constrained by a finite template set but is subject to computational limits.

% ═════════════════════════════════════════════════════════════════════════
% CHAPTER 3: State of the Art
% ═════════════════════════════════════════════════════════════════════════
\section{State of the Art}
\label{sec:sota}

\subsection{Signature-Based Scanners}

\textbf{Nuclei} uses YAML templates for CVE matching. \textbf{Nessus} provides plugin-based detection. Both are deterministic and reproducible.

\textbf{Critical Limitation:} Detection is bounded by template library coverage. Vulnerabilities without CVE identifiers or custom implementations cannot be detected.

\subsection{LLM-Based Security Testing}

\textbf{PentestGPT} demonstrates feasibility of LLM-driven penetration testing but lacks structured state management.

\textbf{Happe et al.} show LLMs can identify vulnerabilities in code but do not integrate with runtime testing.

\textbf{Gap:} No production-scale context window management. No structured hypothesis tracking. No explicit deduplication mechanisms.

\subsection{Agent Architectures}

\textbf{ReAct} (Reasoning + Acting) provides a loop structure: observe $\rightarrow$ reason $\rightarrow$ act.

\textbf{Reflexion} adds self-evaluation but requires additional LLM calls.

\textbf{Limitation:} Most architectures do not address long-context degradation or memory compression.

% ═════════════════════════════════════════════════════════════════════════
% CHAPTER 4: Identified Gaps
% ═════════════════════════════════════════════════════════════════════════
\section{Identified Gaps}
\label{sec:gaps}

\begin{table}[H]
\centering
\caption{Research Gaps Addressed by \phantom{}}
\label{tab:gaps}
\begin{tabularx}{\textwidth}{@{}l X l@{}}
\toprule
\textbf{Gap} & \textbf{Description} & \textbf{Addressed By} \\
\midrule
G1 & No production-scale context window management & Dual-anchor memory (Sec.~\ref{sec:architecture}) \\
G2 & No explicit hypothesis deduplication & Exact-match ledger (Sec.~\ref{sec:formal}) \\
G3 & CVE Template Gap unnamed & Formal definition (Sec.~\ref{sec:problem}) \\
G4 & No bounded termination guarantees & Explicit iteration limits (Sec.~\ref{sec:design}) \\
G5 & No structured state sharing & Shared ledger architecture (Sec.~\ref{sec:architecture}) \\
\bottomrule
\end{tabularx}
\end{table}

% ═════════════════════════════════════════════════════════════════════════
% CHAPTER 5: Proposed Approach
% ═════════════════════════════════════════════════════════════════════════
\section{Proposed Approach}
\label{sec:approach}

\subsection{Design Principles}

\begin{enumerate}[label=\textbf{P\arabic*:}, leftmargin=2em]
    \item \textbf{Bounded Execution:} All loops have explicit termination conditions.
    \item \textbf{Structured State:} Vulnerability hypotheses tracked explicitly, not inferred from conversation.
    \item \textbf{Exact Deduplication:} No semantic similarity; equality-based deduplication for predictability.
    \item \textbf{Dual-Anchor Memory:} Two independent mechanisms preserve critical findings.
    \item \textbf{Optional Sandboxing:} Execution isolation available but not mandatory.
\end{enumerate}

\subsection{System Overview}

\phantom{} implements a ReAct loop with:
\begin{itemize}[leftmargin=*]
    \item Iteration limit (default: 300)
    \item HypothesisLedger for deduplication
    \item MemoryCompressor with anchor extraction
    \item Optional Docker sandboxing
\end{itemize}

% ═════════════════════════════════════════════════════════════════════════
% CHAPTER 6: Formal Model
% ═════════════════════════════════════════════════════════════════════════
\section{Formal Model}
\label{sec:formal}

\subsection{CVE Template Gap}

\begin{definition}[CVE Template Gap]
Let $T$ be a finite template set. Let $V$ be the vulnerability space. The CVE Template Gap is:
\begin{equation}
\forall v \in V : v \notin T \implies \neg\text{detectable}_T(v)
\end{equation}
\end{definition}

\begin{theorem}[Structural Impossibility]
For any signature-based scanner with template set $T$, there exist vulnerabilities $v \in V$ such that $v$ is undetectable.
\end{theorem}

\begin{proof}
Since $T$ is finite and $V$ is unbounded (new vulnerabilities discovered continuously), $|V| > |T|$. Therefore, $\exists v \in V : v \notin T$. By Definition 4.1, such $v$ is undetectable.
\end{proof}

\subsection{HypothesisLedger Formalization}

\begin{definition}[Hypothesis Space]
Let $H \subseteq S \times C$ where:
\begin{itemize}[leftmargin=*]
    \item $S$ = attack surface strings
    \item $C$ = vulnerability class strings
\end{itemize}
\end{definition}

\begin{definition}[Membership Function]
\begin{equation}
\mu((s, c), H) = \begin{cases} 1 & \text{if } \exists (s', c') \in H : s = s' \land c = c' \\ 0 & \text{otherwise} \end{cases}
\end{equation}
\end{definition}

\begin{property}[Exact-Match Deduplication]
The add operation:
\begin{equation}
\text{add}(s, c) = \begin{cases} h'.\text{id} & \text{if } \mu((s, c), H) = 1 \land h' = \text{find}(s, c) \\ \text{new\_id} & \text{otherwise} \end{cases}
\end{equation}
\end{property}

\begin{remark}[Limitation]
Exact-match deduplication treats $(/api/user/1, \text{sqli})$ and $(/api/user/2, \text{sqli})$ as distinct. No semantic equivalence is computed.
\end{remark}

\subsection{Dual-Anchor Memory}

\begin{definition}[Memory System]
Let:
\begin{itemize}[leftmargin=*]
    \item $M$ = message history (ordered list)
    \item $A$ = anchor set (high-signal items)
    \item $C(M)$ = compression function
\end{itemize}
\end{definition}

\begin{property}[Anchor Invariant]
\begin{equation}
A \subseteq C(M) \cup A_s
\end{equation}
where $A_s$ = state-level anchors (\code{finding\_anchors}).
\end{property}

\begin{definition}[Anchor Types]
\begin{enumerate}[label=(\alph*)]
    \item \textbf{Message-level:} Extracted via keyword matching before compression
    \item \textbf{State-level:} Stored separately from conversation history
\end{enumerate}
\end{definition}

\subsection{ReAct Loop Formalization}

\begin{definition}[State Vector]
At iteration $t$:
\begin{equation}
S_t = (M_t, A_t, H_t, C_t, I_t)
\end{equation}
where:
\begin{itemize}[leftmargin=*]
    \item $M_t$ = message history
    \item $A_t$ = anchor set
    \item $H_t$ = hypothesis ledger state
    \item $C_t$ = coverage tracker state
    \item $I_t$ = iteration counter
\end{itemize}
\end{definition}

\begin{definition}[Transition Function]
\begin{equation}
\delta : S_t \times \text{Action} \rightarrow S_{t+1}
\end{equation}
\end{definition}

\begin{algorithm}[H]
\caption{Bounded ReAct Loop}
\label{alg:react}
\begin{algorithmic}[1]
\STATE $t \leftarrow 0$
\WHILE{$\neg\text{should\_stop}(S_t)$}
    \STATE $O_t \leftarrow \text{observe}(S_t)$
    \STATE $T_t \leftarrow \text{reason}(O_t)$
    \STATE $A_t \leftarrow \text{act}(T_t)$
    \STATE $S_{t+1} \leftarrow \text{update}(S_t, A_t, O_t)$
    \STATE $t \leftarrow t + 1$
\ENDWHILE
\STATE \textbf{return} $S_t$
\end{algorithmic}
\end{algorithm}

\begin{definition}[Termination Conditions]
\begin{equation}
\text{should\_stop}(S_t) = \text{completed} \lor (I_t \geq I_{\text{max}}) \lor (n_{\text{no\_action}} \geq 8)
\end{equation}
\end{definition}

% ═════════════════════════════════════════════════════════════════════════
% CHAPTER 7: System Architecture
% ═════════════════════════════════════════════════════════════════════════
\section{System Architecture}
\label{sec:architecture}

\subsection{Component Overview}

\begin{table}[H]
\centering
\caption{Core Components}
\label{tab:components}
\begin{tabularx}{\textwidth}{@{}l l X@{}}
\toprule
\textbf{Component} & \textbf{Location} & \textbf{Function} \\
\midrule
BaseAgent & \code{agents/base\_agent.py} & Abstract ReAct loop implementation \\
PhantomAgent & \code{agents/PhantomAgent/} & Concrete penetration testing agent \\
AgentState & \code{agents/state.py} & Core state storage \\
HypothesisLedger & \code{agents/hypothesis\_ledger.py} & Hypothesis tracking with deduplication \\
CoverageTracker & \code{agents/coverage\_tracker.py} & Surface coverage recording \\
MemoryCompressor & \code{llm/memory\_compressor.py} & Context window management \\
ToolExecutor & \code{tools/executor.py} & Tool validation and execution \\
\bottomrule
\end{tabularx}
\end{table}

\subsection{Execution Flow}

\begin{figure}[H]
\centering
\begin{tikzpicture}[
    node distance=0.8cm and 1.2cm,
    box/.style={rectangle, draw, minimum width=3.5cm, minimum height=0.8cm, align=center},
    decision/.style={diamond, draw, aspect=2, minimum width=2cm, align=center},
    arrow/.style={->, >=Stealth}
]
\node[box] (init) {Initialize State};
\node[decision, below=of init] (stop) {Stop?};
\node[box, right=2.5cm of stop] (done) {Return Results};
\node[box, below=of stop] (observe) {Observe};
\node[box, below=of observe] (reason) {Reason (LLM)};
\node[box, below=of reason] (act) {Act (Tools)};
\node[box, below=of act] (update) {Update State};

\draw[arrow] (init) -- (stop);
\draw[arrow] (stop) -- node[above] {yes} (done);
\draw[arrow] (stop) -- node[left] {no} (observe);
\draw[arrow] (observe) -- (reason);
\draw[arrow] (reason) -- (act);
\draw[arrow] (act) -- (update);
\draw[arrow] (update) -- ++(-3,0) |- (stop);
\end{tikzpicture}
\caption{ReAct Execution Loop}
\label{fig:react}
\end{figure}

\subsection{Data Flow}

\textbf{State $\rightarrow$ LLM:}
\begin{enumerate}[leftmargin=*]
    \item Message history retrieved from \code{AgentState}
    \item MemoryCompressor invoked (threshold-based)
    \item Finding anchors re-injected as user message
    \item Periodic summaries injected (every 10 iterations)
\end{enumerate}

\textbf{LLM $\rightarrow$ Tools:}
\begin{enumerate}[leftmargin=*]
    \item Response parsed for XML tool invocations
    \item Tools validated against registry
    \item Executed in optional sandbox or locally
\end{enumerate}

\textbf{Tools $\rightarrow$ State:}
\begin{enumerate}[leftmargin=*]
    \item Results added as observations
    \item Checkpoint saved every 5 iterations (root agent)
\end{enumerate}

\subsection{Multi-Agent Topology}

Sub-agents share:
\begin{itemize}[leftmargin=*]
    \item \code{HypothesisLedger} (deduplication)
    \item \code{CoverageTracker} (coverage)
\end{itemize}

Each sub-agent owns:
\begin{itemize}[leftmargin=*]
    \item Independent LLM instance
    \item Independent \code{AgentState}
\end{itemize}

% ═════════════════════════════════════════════════════════════════════════
% CHAPTER 8: Design Decisions
% ═════════════════════════════════════════════════════════════════════════
\section{Design Decisions}
\label{sec:design}

\subsection{Bounded Iterations}

\textbf{Decision:} Hard limit of 300 iterations per scan.

\textbf{Rationale:} Prevents infinite loops and unbounded costs.

\textbf{Warning Gates:} Notifications at 80\%, 85\%, and max-3 iterations.

\subsection{Exact-Match Deduplication}

\textbf{Decision:} Deduplication by $(surface, vuln\_class)$ string equality.

\textbf{Rationale:} Predictable behavior; no ambiguity in equivalence.

\textbf{Trade-off:} Syntactic variants treated as distinct hypotheses.

\subsection{Dual-Anchor Memory}

\textbf{Decision:} Two independent anchor mechanisms.

\textbf{Rationale:} Defense in depth against compression loss.

\textbf{Mechanisms:}
\begin{enumerate}[leftmargin=*]
    \item Message-level: Regex keyword matching (50+ patterns)
    \item State-level: \code{finding\_anchors} list
\end{enumerate}

\subsection{Optional Sandboxing}

\textbf{Decision:} Docker sandbox available but not mandatory.

\textbf{Rationale:} Flexibility for development vs. production.

\textbf{Default:} Local execution (\code{PHANTOM\_SANDBOX\_MODE=false}).

\subsection{Checkpoint Frequency}

\textbf{Decision:} Atomic checkpoint every 5 iterations.

\textbf{Rationale:} Limits maximum work loss to 5 iterations.

\textbf{Protocol:} JSON serialization $\rightarrow$ HMAC-SHA256 $\rightarrow$ atomic rename.

% ═════════════════════════════════════════════════════════════════════════
% CHAPTER 9: Limitations and Failure Modes
% ═════════════════════════════════════════════════════════════════════════
\section{Limitations and Failure Modes}
\label{sec:limitations}

\subsection{Memory Limitations}

\begin{property}[Anchor Keyword Dependence]
Anchors extracted via 50+ hardcoded keywords. Findings without keyword matches are lost during compression.
\end{property}

\begin{property}[Compression Lossiness]
LLM-based summarization paraphrases content. Technical details may be generalized despite preservation prompts.
\end{property}

\subsection{Reasoning Limitations}

\begin{property}[Agent Inconsistency]
Sub-agents hold independent contexts. No consensus protocol ensures agreement on hypothesis status.
\end{property}

\begin{property}[No Shared Reasoning]
\code{HypothesisLedger} stores factual state only. Reasoning chains and decision rationales are not shared.
\end{property}

\subsection{Deduplication Limitations}

\begin{property}[Exact-Match Weakness]
$(/api/user/1, \text{sqli})$ and $(/api/user/2, \text{sqli})$ are distinct hypotheses. No semantic equivalence computed.
\end{property}

\subsection{Execution Limitations}

\begin{property}[Tool Hallucination]
No verification that executed tool matches LLM intent. Parameters executed literally.
\end{property}

\begin{property}[RBAC Dependence]
Access control requires RBAC module presence. Lazy import provides no failure-closed default.
\end{property}

\subsection{Boundary Conditions}

\begin{table}[H]
\centering
\caption{System Boundaries}
\label{tab:boundaries}
\begin{tabularx}{\textwidth}{@{}l X@{}}
\toprule
\textbf{Parameter} & \textbf{Bound} \\
\midrule
Max iterations & 300 (default), configurable \\
Checkpoint interval & 5 iterations \\
Context ceiling & 80,000 tokens (regardless of model capacity) \\
No-action abort & 8 consecutive iterations \\
Rate limit retries & 10 consecutive failures \\
Compression chunk size & 10 messages (configurable) \\
\bottomrule
\end{tabularx}
\end{table}

% ═════════════════════════════════════════════════════════════════════════
% CHAPTER 10: Discussion
% ═════════════════════════════════════════════════════════════════════════
\section{Discussion}
\label{sec:discussion}

\subsection{Evaluation Status}

\textbf{Current State:} System implementation verified against codebase. No empirical evaluation has been performed.

\textbf{Proposed Methodology:}
\begin{enumerate}[leftmargin=*]
    \item \textbf{Benchmark:} OWASP Juice Shop, DVWA, WebGoat (intentionally vulnerable applications)
    \item \textbf{Metrics:} Precision, recall, time-to-detection, token consumption
    \item \textbf{Baseline:} Nuclei (signature-based), manual testing
    \item \textbf{Variants:} With/without sandbox, different LLM providers
\end{enumerate}

\subsection{Comparison to Signature-Based Tools}

\begin{table}[H]
\centering
\caption{Paradigm Comparison}
\label{tab:comparison}
\begin{tabularx}{\textwidth}{@{}l X X@{}}
\toprule
\textbf{Property} & \textbf{Signature-Based} & \textbf{LLM-Based} \\
\midrule
Determinism & High & Low (sampling-dependent) \\
Reproducibility & Exact & Approximate \\
Coverage bound & Template library & Reasoning capability \\
False positive rate & Low & Higher (unverified) \\
Computational cost & Low & High (LLM API calls) \\
Explainability & Template match & Reasoning trace \\
\bottomrule
\end{tabularx}
\end{table}

\subsection{Positioning}

\phantom{} occupies a design point between:
\begin{itemize}[leftmargin=*]
    \item Signature-based scanners (deterministic, bounded coverage)
    \item Manual penetration testing (unbounded coverage, unscalable)
\end{itemize}

The system trades determinism for flexibility, bounded only by LLM reasoning capabilities and iteration limits.

% ═════════════════════════════════════════════════════════════════════════
% CHAPTER 11: Conclusion
% ═════════════════════════════════════════════════════════════════════════
\section{Conclusion}
\label{sec:conclusion}

\subsection{Summary}

This thesis presented \phantom{}, an LLM-driven vulnerability detection system with:
\begin{enumerate}[leftmargin=*]
    \item Formal definition of the CVE Template Gap
    \item Bounded ReAct execution with explicit termination
    \item Dual-anchor memory architecture
    \item Exact-match hypothesis deduplication
\end{enumerate}

\subsection{Contributions}

\begin{enumerate}[leftmargin=*]
    \item \textbf{Theoretical:} CVE Template Gap as structural impossibility theorem
    \item \textbf{Architectural:} Dual-anchor memory with verified implementation
    \item \textbf{Formal:} Exact-match deduplication semantics
    \item \textbf{Practical:} Bounded execution guarantees
\end{enumerate}

\subsection{Future Work}

\begin{enumerate}[leftmargin=*]
    \item Empirical evaluation on vulnerability benchmarks
    \item Semantic deduplication exploration
    \item Cross-agent consensus protocols
    \item Reasoning chain preservation mechanisms
\end{enumerate}

% ═════════════════════════════════════════════════════════════════════════
% End Document
% ═════════════════════════════════════════════════════════════════════════
\end{document}
